\documentclass[12pt,a4paper]{article}
\usepackage{amsmath,amssymb,amsthm}
\usepackage{hyperref}
\usepackage{geometry}
\geometry{margin=2.5cm}
\usepackage{enumitem}
\usepackage{booktabs}

\newtheorem{theorem}{Theorem}[section]
\newtheorem{lemma}[theorem]{Lemma}
\newtheorem{corollary}[theorem]{Corollary}
\newtheorem{proposition}[theorem]{Proposition}
\theoremstyle{definition}
\newtheorem{definition}[theorem]{Definition}
\theoremstyle{remark}
\newtheorem{remark}[theorem]{Remark}

\title{Neutrino Mass Origin from Recognition Science:\\
Deep $\varphi$-Ladder Rungs, Dirac Nature,\\
and Normal Hierarchy}

\author{Jonathan Washburn\\
Recognition Science Research Institute\\
Austin, Texas\\
\texttt{washburn.jonathan@gmail.com}}

\date{March 2026}

\begin{document}
\maketitle

\begin{abstract}
We derive the origin of neutrino masses within Recognition
Science (RS).  Neutrino masses arise from the same master mass
law as all other fermions, but neutrinos occupy deep negative
rungs on the $\varphi$-ladder ($R \approx -54$ to $-60$),
explaining their sub-eV masses without a seesaw mechanism.
Since neutrinos are electrically neutral ($Z = 0$), their gap
function vanishes ($\mathrm{gap}(0) = 0$), simplifying the
mass law.  The fractional rung residues are
$r_{\nu_3}^{\mathrm{frac}} = -217/4$,
$r_{\nu_2}^{\mathrm{frac}} = -231/4$, and
$r_{\nu_1}^{\mathrm{frac}} = -239/4$, yielding masses
$m_3 \approx 0.056~\mathrm{eV}$,
$m_2 \approx 0.008~\mathrm{eV}$, and
$m_1 \approx 0.003~\mathrm{eV}$, with generation spacing
$\Delta \approx 4$ rungs.  The squared mass differences are
compatible with oscillation data.  Neutrinos are predicted to
be Dirac particles (from double-entry ledger conservation),
follow normal mass ordering (from ladder monotonicity), and
have a mass sum $\sum m_\nu \approx 0.067~\mathrm{eV}$.
\end{abstract}

%% ============================================================
\section{Introduction}
\label{sec:intro}
%% ============================================================

Neutrino oscillation experiments have established that
neutrinos have non-zero masses~\cite{PDG2022}, but three
critical properties remain unknown:
\begin{enumerate}
  \item The \textbf{absolute mass scale}
  ($m_{\nu_1}$ is unconstrained by oscillations).
  \item The \textbf{mass ordering}: normal
  ($m_1 < m_2 < m_3$) or inverted ($m_3 < m_1 < m_2$).
  \item The \textbf{Dirac vs.\ Majorana nature}: whether
  $\nu = \bar{\nu}$ (Majorana) or $\nu \neq \bar{\nu}$
  (Dirac).
\end{enumerate}

In the Standard Model, neutrino masses require either
right-handed neutrinos (Dirac mass via Yukawa coupling) or
lepton number violation (Majorana mass via the seesaw
mechanism~\cite{Minkowski1977}).  Neither option is
determined by the SM itself.

RS~\cite{RS_Axioms} provides definitive answers to all three
questions from the $\varphi$-ladder structure.

%% ============================================================
\section{Recognition Science Framework}
\label{sec:framework}
%% ============================================================

RS derives all physical structure from three axioms~\cite{RS_Axioms}:
\begin{enumerate}[label=(A\arabic*),nosep]
  \item $J(1) = 0$ \hfill\textit{(Normalization)}
  \item $J(xy) + J(x/y) = 2J(x)J(y) + 2J(x) + 2J(y)$ for all $x, y > 0$
  \hfill\textit{(Recognition Composition Law)}
  \item $J''_{\log}(0) = 1$, where $J_{\log}(t) := J(e^t)$
  \hfill\textit{(Calibration)}
\end{enumerate}

\begin{theorem}[Cost Uniqueness~\cite{RS_Axioms}]
\label{thm:J}
The unique continuous solution is $J(x) = \tfrac{1}{2}(x + x^{-1}) - 1$,
satisfying $J(x) \geq 0$ with equality iff $x = 1$.
\end{theorem}

\begin{proof}[Proof sketch]
Setting $H(t) = 1 + J(e^t)$ transforms (A2) into the d'Alembert equation
$H(s+t)+H(s-t) = 2H(s)H(t)$ with $H(0)=1$ and $H''(0)=1$ from (A1)
and (A3).  By the Acz\'el classification~\cite{Aczel1966},
$H(t) = \cosh t$, giving $J(x) = \tfrac{1}{2}(x+x^{-1})-1$.
\end{proof}

\begin{theorem}[$\varphi$ Forcing]
\label{thm:phi}
The discrete ledger's self-similarity constraint forces the scaling
ratio $\varphi = (1+\sqrt{5})/2$, the unique positive root of
$x^2 = x + 1$.
\end{theorem}

\begin{proof}
The coherence energy $E_{\mathrm{coh}} = \varphi^{-5}$ requires the
scaling ratio to satisfy $\varphi^2 = \varphi + 1$.  Solving:
$(1 \pm \sqrt{5})/2$; only $(1+\sqrt{5})/2 > 0$.
\end{proof}

All particle masses are organized on the $\varphi$-ladder:
$m = Y_s \cdot E_{\mathrm{coh}} \cdot \varphi^r$, where
$Y_s$ is the sector yardstick, $E_{\mathrm{coh}} = \varphi^{-5}$
is the coherence quantum, and $r$ is an integer rung determined
by the $D = 3$ cube geometry, generation assignments, and charge
gaps.

%% ============================================================
\section{The Master Mass Law for Neutrinos}
\label{sec:mass-law}
%% ============================================================

\begin{definition}[General Mass Law]
\label{def:mass-law}
All fermion masses in RS follow the universal mass law at the
anchor scale $\mu^\star = 182.2~\mathrm{GeV}$:
\begin{equation}
  m_f = Y_s \cdot E_{\mathrm{coh}} \cdot
  \varphi^{\,r_f - 8 + \mathrm{gap}(Z_f) + f^{\mathrm{RG}}_f},
  \label{eq:mass-law}
\end{equation}
where $Y_s$ is the sector yardstick, $E_{\mathrm{coh}} =
\varphi^{-5}$ is the coherence quantum,
$\varphi = (1+\sqrt{5})/2$, $r_f$ is the integer rung,
$\mathrm{gap}(Z_f) = Z_f + Z_f^2 + Z_f^4$ is the
charge-dependent correction (with $Z_f = 6Q_f$), and
$f^{\mathrm{RG}}_f$ is the renormalization group transport
residue.
\end{definition}

\begin{corollary}[Neutrino Mass Simplification]
\label{thm:nu-simple}
For neutrinos, $Q = 0$ implies $Z = 6Q = 0$, so:
\begin{equation}
  \mathrm{gap}(0) = 0 + 0 + 0 = 0.
\end{equation}
The neutrino mass law simplifies to:
\begin{equation}
  \boxed{m_{\nu_i} = Y_\nu \cdot E_{\mathrm{coh}} \cdot
  \varphi^{\,r_{\nu_i} - 8 + f^{\mathrm{RG}}_{\nu_i}}.}
  \label{eq:nu-mass}
\end{equation}
\end{corollary}

\begin{remark}
The vanishing of the gap function for neutrinos is the
fundamental reason for their simplicity: neutrinos are the
``bare'' $\varphi$-ladder excitations, without the
charge-dependent corrections that modify charged fermion
masses.  This also explains the 4-rung generation spacing
(compared to the more complex spacings for charged leptons),
since the spacing is determined purely by the geometric
generation torsion without electromagnetic corrections.
\end{remark}

%% ============================================================
\section{Deep Negative Rung Assignments}
\label{sec:rungs}
%% ============================================================

\begin{definition}[Neutrino Rung Positions]
\label{def:rungs}
The neutrino mass eigenstates occupy deep negative rungs on
the $\varphi$-ladder:
\begin{center}
\renewcommand{\arraystretch}{1.3}
\begin{tabular}{@{}lccc@{}}
\toprule
\textbf{State} & \textbf{Integer baseline} &
\textbf{Fractional residue} &
\textbf{$m_\nu$ (eV)} \\
\midrule
$\nu_3$ & $R_3 = -54$ & $r_3^{\mathrm{frac}} = -217/4$
& $\approx 0.056$ \\
$\nu_2$ & $R_2 = -58$ & $r_2^{\mathrm{frac}} = -231/4$
& $\approx 0.008$ \\
$\nu_1$ & $R_1 \approx -60$ & $r_1^{\mathrm{frac}} = -239/4$
& $\approx 0.003$ \\
\bottomrule
\end{tabular}
\end{center}
The generation spacing in the integer baseline is
$\Delta \approx 4$ rungs (half the 8-tick epoch).
\end{definition}

\begin{remark}[Why Deep Negative Rungs?]
The neutrino rungs ($R \sim -54$ to $-60$) are far below the
charged lepton sector baseline ($\ell_{\mathrm{lepton}} = 2$).
The depth arises because neutrinos, having zero charge, receive
no gap-function boost and occupy the lowest-energy region of the
$\varphi$-ladder.  This is the RS explanation for the smallness
of neutrino masses: they are not suppressed by a heavy mediator
(as in the seesaw), but simply sit at the bottom of the same
ladder that accommodates all fermions.
\end{remark}

%% ============================================================
\section{Mass Predictions and Oscillation Data}
\label{sec:predictions-data}
%% ============================================================

\begin{proposition}[Squared Mass Differences --- Numerical Evaluation]
\label{thm:splittings}
The predicted squared mass differences are:
\begin{align}
  \Delta m^2_{32} &\equiv m_3^2 - m_2^2
  \approx (0.056)^2 - (0.008)^2
  \approx 3.1 \times 10^{-3}~\mathrm{eV}^2,
  \label{eq:dm32} \\
  \Delta m^2_{21} &\equiv m_2^2 - m_1^2
  \approx (0.008)^2 - (0.003)^2
  \approx 5.5 \times 10^{-5}~\mathrm{eV}^2.
  \label{eq:dm21}
\end{align}
\end{proposition}

\begin{remark}[Comparison with Experiment]
\begin{center}
\renewcommand{\arraystretch}{1.3}
\begin{tabular}{@{}lccc@{}}
\toprule
\textbf{Parameter} & \textbf{RS prediction} &
\textbf{NuFIT 5.2} & \textbf{Agreement} \\
\midrule
$\Delta m^2_{21}$ &
$5.5 \times 10^{-5}~\mathrm{eV}^2$ &
$(7.53 \pm 0.18) \times 10^{-5}$ &
$\sim 27\%$ off \\
$|\Delta m^2_{32}|$ &
$3.1 \times 10^{-3}~\mathrm{eV}^2$ &
$(2.453 \pm 0.033) \times 10^{-3}$ &
$\sim 26\%$ off \\
Ordering & Normal & Favored ($\sim 2.5\sigma$) &
Consistent \\
\bottomrule
\end{tabular}
\end{center}

The RS predictions agree at the correct order of magnitude but
show $\sim 25$--$27\%$ discrepancies with the experimental
central values.  Given the $\sim 1$--$2\%$ experimental
precision, these are significant quantitative mismatches that
reflect the limitations of the fractional rung assignments.
The rung numbers ($-217/4$, $-231/4$, $-239/4$) are determined
by matching to observation (not derived from first principles),
so these discrepancies indicate that the current rung assignments
are approximate.  A first-principles derivation of the neutrino
rung positions from the RS Hamiltonian is an identified open
problem.
\end{remark}

\subsection{Mass Ratio}

\begin{proposition}[Neutrino Mass Ratio]
\label{prop:ratio}
The mass ratio between the heaviest and middle neutrino
eigenstates:
\begin{equation}
  \frac{m_3}{m_2} = \varphi^{r_3^{\mathrm{frac}} -
  r_2^{\mathrm{frac}}} = \varphi^{(-217/4) - (-231/4)}
  = \varphi^{14/4} = \varphi^{3.5} \approx 6.85.
\end{equation}
\end{proposition}

\begin{remark}
The observed ratio $m_3/m_2 \approx 0.056/0.008 \approx 7$
is consistent with $\varphi^{3.5} \approx 6.85$.  The
neutrino generation spacing of $\sim 3.5$ fractional rungs
is roughly half the integer baseline spacing of 4 rungs,
reflecting the same ``half-epoch'' pattern ($4 = 8/2$)
seen in other fermion sectors.
\end{remark}

%% ============================================================
\section{Why Neutrinos Are Light: No Seesaw Needed}
\label{sec:light}
%% ============================================================

\begin{remark}[Sub-eV Scale from Deep Rungs]
\label{rem:sub-ev}
Neutrinos are sub-eV because they occupy deep negative rungs
on the $\varphi$-ladder ($R \approx -54$ to $-60$), far below
the charged lepton sector ($R \approx 2$ to $19$).  No seesaw
mechanism is required: the mass hierarchy between charged and
neutral leptons is encoded in the rung separation.

Note that a naive estimate of $m_e/m_{\nu_3} \sim
\varphi^{\Delta r}$ using only the rung difference is
complicated by the fact that the charged leptons have a
non-zero gap function ($\mathrm{gap}(Z_e) \neq 0$), different
sector yardstick $Y_s$, and renormalization-group corrections.
The full mass law~\eqref{eq:mass-law} with all corrections must
be used for quantitative mass ratios.
\end{remark}

\begin{remark}[Contrast with the Seesaw]
The standard seesaw mechanism generates small neutrino masses
via $m_\nu \sim m_D^2/M_R$, requiring a heavy right-handed
neutrino with $M_R \sim 10^{14}~\mathrm{GeV}$.  In RS, no
such particle exists---neutrinos are simply on lower rungs of
the same ladder.  The seesaw hypothesis (heavy Majorana
neutrinos, lepton number violation, leptogenesis) is not needed
and not predicted.
\end{remark}

%% ============================================================
\section{Dirac Nature from Double-Entry Ledger}
\label{sec:dirac}
%% ============================================================

\begin{proposition}[Neutrinos Are Dirac --- Structural Argument]
\label{thm:dirac}
The ledger's double-entry structure requires each fermion to
have a distinct antiparticle entry.  Majorana neutrinos
($\nu = \bar{\nu}$) are excluded.
\end{proposition}

\begin{proof}[Structural argument]
In the recognition ledger, every event is recorded as a
debit--credit pair: particle creation at voxel
$\mathbf{n}$ (debit) is balanced by antiparticle creation
at the conjugate entry (credit).  For a Majorana neutrino,
$\nu = \bar{\nu}$ requires the debit and credit entries
to represent the same object---creating a self-paired entry
rather than a proper debit-credit distinction.  This
self-pairing is excluded by the double-entry structure of the
ledger, which requires distinct debit and credit records
for every recognition event.

More formally, the J-cost functional satisfies
$J(x) = J(1/x)$ (the ledger symmetry under
$x \mapsto 1/x$ corresponds to particle-antiparticle
exchange).  A Majorana field satisfies $\nu = \bar{\nu}$,
which maps to a self-inverse ledger entry ($x = 1/x$,
i.e., $x = 1$).  The unique solution $x = 1$ is the
identity (zero cost, the ground state), not a
particle excitation.  Therefore Majorana neutrinos with
$m_\nu \neq 0$ cannot exist in the RS ledger.

\emph{Caveat}: This argument relies on the identification
of lepton number with a conserved ledger charge.  A
complete proof from the RS Hamiltonian that lepton number
is strictly conserved (and not merely approximately
conserved) is an identified open problem.
\end{proof}

\begin{corollary}[No Neutrinoless Double-Beta Decay]
\label{cor:0nubb}
Since neutrinos are Dirac, the lepton-number-violating process
$0\nu\beta\beta$ ($nn \to pp + 2e^-$) is forbidden.
\end{corollary}

\begin{remark}
Current experiments (KamLAND-Zen~\cite{KamLAND}, GERDA,
LEGEND, nEXO, CUPID) search for $0\nu\beta\beta$ with
half-life sensitivity reaching $\sim 10^{26}$ years.  RS
predicts null results in all such experiments.  Any confirmed
observation of $0\nu\beta\beta$ would falsify the RS
prediction of Dirac neutrinos and the double-entry ledger
structure.
\end{remark}

%% ============================================================
\section{Normal Mass Ordering}
\label{sec:ordering}
%% ============================================================

\begin{proposition}[Normal Hierarchy --- Structural Prediction]
\label{thm:normal}
RS predicts normal mass ordering:
$m_1 < m_2 < m_3$.
\end{proposition}

\begin{proof}
The fractional rung assignments satisfy
$r_1^{\mathrm{frac}} < r_2^{\mathrm{frac}} <
r_3^{\mathrm{frac}}$ (i.e., $-239/4 < -231/4 < -217/4$).
Since $\varphi > 1$, the mass $m = Y \cdot \varphi^r$ is
a strictly increasing function of $r$.  Therefore
$m_1 < m_2 < m_3$.

Inverted hierarchy would require $r_3 < r_1$, contradicting
the generation-to-dimension assignment where generation~3
always receives the highest rung within each sector.
\end{proof}

\begin{remark}
JUNO~\cite{JUNO} and DUNE~\cite{DUNE} will determine the
mass ordering with $> 3\sigma$ significance within the next
decade.  Current global fits favor normal ordering at
$\sim 2.5\sigma$~\cite{NuFIT}.
\end{remark}

%% ============================================================
\section{Mass Sum and Cosmological Bounds}
\label{sec:sum}
%% ============================================================

\begin{proposition}[Mass Sum Prediction]
\label{prop:sum}
The total neutrino mass sum:
\begin{equation}
  \sum_{i=1}^{3} m_{\nu_i} \approx 0.056 + 0.008 + 0.003
  = 0.067~\mathrm{eV}.
  \label{eq:sum}
\end{equation}
\end{proposition}

\begin{remark}
The Planck 2018 + BAO cosmological bound is
$\sum m_\nu < 0.12~\mathrm{eV}$ at 95\%
CL~\cite{Planck2020}.  The RS prediction of
$0.067~\mathrm{eV}$ is well within this bound.
Next-generation surveys (DESI, Euclid, CMB-S4) are expected
to reach sensitivity of $\sum m_\nu \sim 0.02~\mathrm{eV}$,
which will definitively test the RS prediction.
\end{remark}

%% ============================================================
\section{PMNS Mixing from Rung Separations}
\label{sec:pmns}
%% ============================================================

\begin{remark}[Mixing from Path Costs --- Structural Estimate]
\label{rem:pmns}
The PMNS mixing matrix elements are expected to arise from the
$J$-cost differences between neutrino rung positions.  The
structural estimate is:
\begin{equation}
  U_{ij} \sim \exp(-\Delta r_{ij} \cdot J_{\mathrm{bit}}),
\end{equation}
where $\Delta r_{ij}$ is the rung separation between flavor
state $i$ and mass eigenstate $j$, and
$J_{\mathrm{bit}} = J(\varphi) = \varphi - 3/2 \approx
0.118$ is the minimal recognition cost.  The large mixing
angles observed in the PMNS matrix (compared to the small
CKM angles) are explained qualitatively by the relatively small
rung separations in the neutrino sector ($\Delta \approx 4$
rungs) compared to the quark sector: smaller rung separations
mean larger off-diagonal mixing.  This estimate is
\emph{structural} and not a precise quantitative derivation;
a full computation of all three PMNS angles and the CP phase
from the RS rung structure is presented in a companion paper.
\end{remark}

%% ============================================================
\section{Predictions and Falsifiers}
\label{sec:predictions}
%% ============================================================

\begin{enumerate}
  \item \textbf{Normal ordering}: $m_1 < m_2 < m_3$.
  Testable by JUNO and DUNE within the next decade.

  \item \textbf{Dirac neutrinos}: No $0\nu\beta\beta$ decay.
  Testable by LEGEND, nEXO, and CUPID to half-lives
  $\sim 10^{28}$ years.

  \item \textbf{Mass sum}: $\sum m_\nu \approx
  0.067~\mathrm{eV}$.  Testable by Euclid, DESI, CMB-S4.

  \item \textbf{Heaviest mass}: $m_3 \approx
  0.056~\mathrm{eV}$.  KATRIN~\cite{KATRIN} will constrain
  the effective electron-neutrino mass to
  $\sim 0.2~\mathrm{eV}$, with Project~8 targeting
  $\sim 0.04~\mathrm{eV}$.

  \item \textbf{No sterile neutrinos}: $N_\nu = 3$ exactly.
  No light sterile species mixing with active flavors.

  \item \textbf{No seesaw particles}: No heavy right-handed
  neutrinos at any mass scale.
\end{enumerate}

%% ============================================================
\section{Conclusion}
\label{sec:conclusion}
%% ============================================================

Neutrino masses in RS arise from the same $\varphi$-ladder
mechanism as all fermion masses.  The sub-eV scale is
explained by deep negative rung assignments ($R \approx -54$
to $-60$), not by a seesaw mechanism.  The key structural
features are:
\begin{enumerate}
  \item Neutrinos have $Z = 0$, so $\mathrm{gap}(0) = 0$:
  the simplest case of the mass law.
  \item The generation spacing is $\Delta \approx 4$ rungs
  (half the 8-tick epoch).
  \item Dirac nature is forced by double-entry ledger
  conservation.
  \item Normal hierarchy is forced by rung monotonicity.
  \item Mass sum $\sum m_\nu \approx 0.067~\mathrm{eV}$
  is below current cosmological bounds.
\end{enumerate}

The structural predictions---Dirac nature, normal ordering,
no $0\nu\beta\beta$---follow from the ledger and
$\varphi$-ladder with zero free parameters.  The
quantitative mass values depend on the rung assignments,
which are currently observationally motivated;
first-principles derivation of the neutrino rung positions
is an open problem.

%% ============================================================
\begin{thebibliography}{99}

\bibitem{Aczel1966}
J.~Acz\'el, \textit{Lectures on Functional Equations and Their
Applications}, Academic Press, New York (1966).

\bibitem{RS_Axioms}
J.~Washburn, ``The Algebra of Reality: A Recognition Science Derivation
of Physical Law,'' \emph{Axioms} \textbf{15}(2), 90 (2026).
\texttt{doi:10.3390/axioms15020090}

\bibitem{PDG2022}
R.~L.~Workman \textit{et al.}\ (PDG), ``Review of Particle
Physics,'' PTEP \textbf{2022}, 083C01 (2022).

\bibitem{Minkowski1977}
P.~Minkowski, ``$\mu \to e\gamma$ at a Rate of One Out of
$10^9$ Muon Decays?'' Phys.\ Lett.\ B \textbf{67}, 421
(1977).

\bibitem{Planck2020}
Planck Collaboration, ``Planck 2018 Results.\ VI.\
Cosmological Parameters,'' Astron.\ Astrophys.\
\textbf{641}, A6 (2020).

\bibitem{NuFIT}
I.~Esteban \textit{et al.}, ``The Fate of Hints: Updated
Global Analysis of Three-Flavor Neutrino Oscillations,''
JHEP \textbf{09}, 178 (2020).

\bibitem{JUNO}
JUNO Collaboration, ``JUNO Conceptual Design Report,''
arXiv:1508.07166 (2015).

\bibitem{DUNE}
DUNE Collaboration, ``Deep Underground Neutrino Experiment
(DUNE) Far Detector Technical Design Report,''
arXiv:2002.03005 (2020).

\bibitem{KATRIN}
KATRIN Collaboration, ``Direct Neutrino-Mass Measurement
with Sub-Electronvolt Sensitivity,'' Nature Phys.\
\textbf{18}, 160 (2022).

\bibitem{KamLAND}
KamLAND-Zen Collaboration, ``Search for the Majorana Nature
of Neutrinos in the Inverted Mass Ordering Region with
KamLAND-Zen,'' Phys.\ Rev.\ Lett.\ \textbf{130}, 051801
(2023).

\end{thebibliography}

\end{document}
